% ===== Phụ lục C =====
\subsection{Nhãn chuẩn cấu trúc bảng (K0)}
Thí nghiệm K0 (\path{k0_gt.py}) dùng 55 ảnh của \bo{2} có nhãn thân bảng, với \num{3255} cụm từ.
Hai ảnh còn lại chỉ có nhãn loại phiếu.
Nhãn dòng được tạo bằng thuật toán hợp-tìm (union-find) trên các cặp cụm từ chồng nhau theo chiều dọc ít nhất 50\%.
Cách tạo nhãn này độc lập với cả hai thuật toán gom dòng được so sánh.
Nhãn cột là vai trò của cụm từ.
Cả 55 ảnh được soát bằng mắt (\path{results/k0/review/}), và 32 ảnh được giữ lại với 498 dòng chỉ số và \num{1953} cụm từ.
Các ảnh bị loại có khung nhãn gốc lệch khỏi chữ, bị chụp mờ, hoặc có bố cục máy in nhiều cột mà nhãn không phân biệt được.

Chín ảnh được sửa nhãn theo quy tắc ghi trong \path{k0_review.json}:
\begin{enumerate}[label=(\roman*)]
  \item bỏ các dòng mô tả (Peripheral Smear, Morphology, ``Adequate'') và các tiêu đề nhóm không có kết quả, tổng cộng 32 cụm từ;
  \item gộp các tên chỉ số viết trên hai dòng và các tên bị lệch xuống dòng dưới do ảnh nghiêng, tổng cộng 3 cụm từ.
\end{enumerate}

\subsection{Nhãn chữ cho thí nghiệm K5}
\Bo{2} không có nhãn chữ.
Chữ trong từng khung do LLM (Claude) phiên âm từ ảnh gốc màu, không dùng ảnh sau xử lý, để nhãn chuẩn không mang lỗi của bước (1) và (2).
Ô không đọc được để trống và bị loại khỏi nhãn chuẩn.
Ô đọc không chắc chắn được đánh dấu để soát trước (\path{k0_text_gt.json}).
LLM độc lập với Tesseract nên không thiên vị phương án nào trong K5.
Trên 32 ảnh giữ lại có \num{1952} khung có chữ, trong đó 23 khung vẫn được đánh dấu là không chắc chắn.

\subsection{Kết quả đầy đủ của K1, K2 và K4}
Bảng~\ref{tab:K2full} trình bày kết quả đầy đủ, gồm cả hai cách đặt tham số theo $h$.
Tỉ lệ ``ảnh đúng hết'' là phần trăm số ảnh không sai cặp cụm từ nào.
Các phương án có ``xóa kẻ'' chạy lược đồ chiếu trên ảnh đã trừ đường kẻ bằng phép mở với SE đường thẳng.

\begin{table}[H]
\centering\small
\caption{Kết quả gom dòng và tìm cột trên 32 ảnh \bo{2}, trung bình $\pm$ độ lệch chuẩn qua 5 fold, với $h$ ước lượng theo từng ảnh hoặc cố định (thí nghiệm K1, K2, K4).
Giá trị tốt nhất trong từng bước con được in đậm.}\label{tab:K2full}
\begin{tabular}{lllrrrr}
\toprule
\textbf{Bước con} & \textbf{Phương án} & \textbf{$h$} & \textbf{F1} & \textbf{Precision} & \textbf{Recall} & \makecell[r]{\textbf{Ảnh đúng}\\\textbf{hết (\%)}} \\
\midrule
Gom dòng & Lược đồ chiếu ngang & Từng ảnh & 0,705 $\pm$ 0,210 & 0,592 & 0,994 & 71,4 \\
         & Lược đồ chiếu ngang & Cố định & 0,814 $\pm$ 0,205 & 0,739 & 0,993 & 74,3 \\
         & \makecell[l]{Lược đồ chiếu ngang,\\xóa kẻ} & Từng ảnh & 0,937 $\pm$ 0,115 & 0,905 & \textbf{0,995} & \textbf{80,5} \\
         & \makecell[l]{Lược đồ chiếu ngang,\\xóa kẻ} & Cố định & 0,942 $\pm$ 0,105 & 0,912 & \textbf{0,995} & \textbf{80,5} \\
         & Tâm $y$, khe cố định & --- & \textbf{0,990 $\pm$ 0,007} & \textbf{1,000} & 0,980 & 64,8 \\
\midrule
Tìm cột  & \makecell[l]{Lược đồ chiếu dọc\\(có hoặc không xóa kẻ)} & \makecell[l]{Từng ảnh\\hoặc cố định} & \textbf{0,963 $\pm$ 0,023} & \textbf{0,930} & \textbf{1,000} & \textbf{87,1} \\
         & XY-cut đệ quy & \makecell[l]{Từng ảnh\\hoặc cố định} & 0,913 $\pm$ 0,057 & 0,847 & 0,997 & 74,3 \\
         & Khe $x$ cố định & --- & 0,893 $\pm$ 0,051 & 0,813 & 0,997 & 65,2 \\
\bottomrule
\end{tabular}
\end{table}

Thí nghiệm K4 so sánh $h$ ước lượng theo từng ảnh với $h$ cố định bằng 5~px, là trung vị của cả bộ.
Với tìm cột, hai cách cho cùng kết quả.
Với gom dòng, $h$ cố định cho F1 cao hơn một chút (\num{0.942} so với \num{0.937} khi đã xóa kẻ).
Nguyên nhân là $h$ ước lượng trên ảnh 640~px chỉ nằm trong khoảng 4--11~px, và mọi ảnh trong \bo{2} có cùng kích thước.
Lợi thế của tham số theo $h$ chỉ xuất hiện khi các ảnh có độ phân giải khác nhau, điều mà \bo{2} không kiểm tra được.

\subsection{Kết quả đầy đủ của K3}
Thí nghiệm K3 (\path{k3_skew.py}, \path{results/k3/table.csv}) cho thấy ước lượng góc của bước (1), dùng lược đồ chiếu với bước $0{,}25^\circ$, có sai số trung vị $0^\circ$ ở mọi góc thử.
F1 của phương án ``chỉnh nghiêng rồi chiếu'' giảm ở góc $4$--$5^\circ$ (Hình~\ref{fig:K3}).
Nguyên nhân chủ yếu là khung cụm từ bị xoay hai lần nên hình chữ nhật bao phình ra, không phải do góc ước lượng sai.

\subsection{Phiếu xét nghiệm Việt Nam chụp thật (K6)}
Thí nghiệm K6 (\path{k6_bo4.py}) kiểm tra định tính 5 phiếu của \bo{4} chọn ngẫu nhiên với seed 42.
Ảnh chứa thông tin bệnh nhân nên không đưa vào báo cáo.
Tìm cột cho kết quả gần đúng ở 2/5 phiếu, khi bảng chiếm gần hết khung hình, nhưng tách thừa cột ở chỗ có nhiều đường kẻ sát nhau.
Gom dòng thất bại ở cả 5 phiếu, vì ba nguyên nhân:
\begin{enumerate}
  \item giấy cong làm đường kẻ không thẳng, nên phép mở với SE đường thẳng không lấy được đường kẻ, ví dụ một phiếu chỉ đúng 8 dòng đầu rồi gộp phần còn lại;
  \item nét khoanh bằng tay cắt ngang nhiều dòng;
  \item 2 phiếu được chụp kèm nền mà bước (1) không tách được trang, nên cả trang bị coi là bảng.
\end{enumerate}
Phương pháp chọn trên \bo{2}, với ảnh 640~px và phần lớn bảng không kẻ, không chuyển được sang phiếu Việt Nam có kẻ và chữ cao 10--27~px.
Để áp dụng cho loại phiếu này, cần tìm vùng bảng trước và lấy đường kẻ bằng cách chịu được đường cong, ví dụ dùng SE ngắn theo từng đoạn hoặc biến đổi Hough, trước khi chiếu.

\subsection{Bảng tham số}
Tham số được chọn bằng tìm kiếm lưới trên 4 fold, lặp lại 5 lượt.
Tham số ở từng fold được lưu trong \path{results/k2/folds.csv}.
Cả 5 lượt chọn cùng một bộ tham số cho mọi phương án đề xuất, cho thấy tham số ổn định.
Bảng~\ref{tab:params-table} liệt kê các tham số.

\begin{table}[H]
\centering\small
\caption{Tham số của bước (3).}\label{tab:params-table}
\begin{tabularx}{\linewidth}{XlX}
\toprule
\textbf{Tham số} & \textbf{Khoảng thử} & \textbf{Giá trị chọn và lý do} \\
\midrule
SE giãn ngang khi gom dòng ($\times h$) & 2, 3, 4, 5 & 2, theo F1 cặp cùng dòng; SE lớn hơn nối sang dòng kề khi các dòng sát nhau \\
Cửa sổ làm trơn $W_s$ & 1, 3, 5 & 1; các dòng chỉ cách nhau vài px nên làm trơn xóa mất thung lũng. Học phần khuyến nghị 3--5 cho histogram mức xám \chuong{5, S62} \\
Ngưỡng thung lũng $T$ (\% đỉnh) & 2, 5, 10 & 10\% \\
Khe tâm $y$ cố định (px) & 3, 5, 8, 10, 15 & 5 (đối chứng) \\
Độ rộng khe cột ($\times h$); tỉ lệ dòng khe xuyên qua & 0,5--2; 40--80\% & 0,5 và 60\%, theo F1 cặp cùng cột \\
Khe XY-cut ($\times h$); khe $x$ cố định (px) & 0,5--2; 5--30 & Theo F1 (đối chứng) \\
SE lấy đường kẻ & $W/20$, $H/20$ & Cố định theo kế hoạch \chuong{7, S6, S28} \\
Hệ số phóng ô trước OCR & $\times3$ & Cố định; chữ \bo{2} cao 4--11~px (trung vị 5~px), trong khi Tesseract cần khoảng 20~px \\
\bottomrule
\end{tabularx}
\end{table}
